% 08_conclusion.tex
\section{Conclusion}
\label{sec:conclusion}

On the periodic windows --- the family that contains paper~E's strips and the combs that undercut them --- the smallest resolution gap
is bounded from both sides with explicit assumptions and computer-assisted constants for every $10^{-100000}\le\sigma\le10^{-24}$
(Theorem~\ref{thm:BF1}), with $\liminf g_{\min}\sqrt{\ln(1/\sigma)}\ge1-\bar\varepsilon$ proved and the limit $1$ derived (Corollary~\ref{cor:limit}). The law paper~E found on the strips is thereby a theorem on a larger family on this range, and its constant is $1$ up to the derived step; the fitted
$1.02$ reflects the slow $O(\ln\ln/\ln)$ approach. At finite blur the family's minimum is the comb, which is the strip seen through a lens
$256$ times sharper; it comes below the best Gaussian mixtures already at $\sigma\approx1.4\times10^{-3}$, not only below $7\times10^{-6}$.

The theorems are asymptotic ($\sigma\le10^{-24}$); a proof of the upper side below $10^{-100000}$ needs explicit remainders of the
leading-order expansion along the comb sequence; at practical resolutions paper~E's results and the numbers of Sec.~\ref{sec:finite} apply.
They are conditional on (P0)--(P7); the a-priori condition of v1.0 on $F_c$ is now a consequence (Lemma~\ref{lem:BJ1}). The most natural next
steps are to extend to $x<1$ (for strips first) and to odd profiles.
For general states, cases~1 and~2 of the case split are proved, and case~3 needs one inequality (Conjecture~\ref{conj:VI}); paper~E's
Conjecture~1 is still open.
